\section{Scott's indexical and intensional framework}
\label{sec:p08-n-b-006-01}

Scott's 1970 \emph{Advice on Modal Logic} develops a semantic framework,
not a modern list of E/EM/EC/EN completeness theorems. Its six sections
distinguish domains of individuals, reference indices, partial individual
concepts, intensional operations, binding operations, and incidence.
The December 1969 postscript is part of the published presentation:
its qualifications matter to the interpretation of the earlier proposals
\cite[pp.~143--173]{Scott1970}. The propositional neighborhood clause
appears inside this larger account on p.~160.
The opening argument distinguishes a translation used to establish
results from a reason to discard the translated subject. In the virtual
object discussion, Scott likewise distinguishes an explicit axiom system
from eliminating each expression separately in its possible contexts:
axioms expose a consistency commitment, while elimination, when available,
is a metatheorem of the chosen theory. These are methodological arguments,
not a supplied general elimination theorem for every virtual-object system.

\subsection{Individuals, virtuals and a fixed reference space}

\begin{hbdefinition}{def:p08-sa-domains}{The three domains and reference indices}
The source fixes a nonempty domain D of possible individuals, a domain V
of virtual individuals with \(D\subseteq V\), and an index set I;
the pointed interpretations compared here use \(I\ne\varnothing\).
Each \(i\in I\) has a possibly empty domain \(A_i\subseteq D\)
of actual individuals. Enlarging D or I changes the interpretation;
it is not an unnoticed variation inside one fixed model. An index may
have several coordinates, such as world, time, position and agent.
Distinct indices need not have distinct actual domains. Thus the indices
are not identified with either \(A_i\) or sets of true sentences.

For each \(a\in V\), the source constant \(\bar a\) has
the rigid value a. Variables and the ordinary quantifiers range over D,
not all V. Source dotted actual quantifiers are written here
\(\forall_{\cdot}\) and \(\exists_{\cdot}\): at index i
they range over \(A_i\). The source \(\widetilde A\) predicate
tests current actuality. These typed source bars, dots and tildes are
not numeral bars or an unannounced normal-modal syntax
\cite[\S\S I--II, pp.~144--151]{Scott1970}.
\end{hbdefinition}

\begin{hbexample}{def:p08-sa-cross-time-domains}{The source's motivation for possible-domain quantification}
Comparisons of officeholders who served at different times, comparisons
of award recipients, and claims about future readers need not concern
individuals who are all actual at the present index. A fixed D allows
the compared individuals to be quantified over, while \(A_i\) specifies
who is currently actual. The source examples motivate the distinction;
they do not prove that every intended English quantifier has this one
formal reading. The number of coordinates and the actual-domain family
must be chosen for the interpretation's particular task.
\end{hbexample}

\begin{hbmethod}{def:p08-sa-virtuals}{Quantification over indexed virtual objects}
In the source's class-theoretic illustration, if sets x have cardinal
classes \(|x|=\{y:y\text{ is equipollent to }x\}\), a cardinal
quantifier is defined by replacing its cardinal variable with \(|x|\)
and quantifying over x. More generally, for a specified map
\(g:D\to V\), define quantification over its image by
\[
 \forall m\,\Phi(m)\defeq\forall x\,\Phi(g(x)),\qquad
 \exists m\,\Phi(m)\defeq\exists x\,\Phi(g(x)).
\]
\(\operatorname{Card}(\tau)\leftrightarrow\exists x(\tau=|x|)\),
with x not free in \(\tau\), defines the source cardinal predicate.
The new m ranges only over \(g[D]\); repeated names of a virtual
do not alter either clause. This is a restricted virtual quantification,
not a silent expansion of all ordinary variables to V. The cardinal
notation here names a virtual equipollence class, not a numerical
cardinality operation on arbitrary neighborhood worlds.
\end{hbmethod}

\begin{hbtheorem}{thm:p08-sa-indexed-virtuals}{proved=annex:p8.sa.indexed-virtuals}{Image quantification and the source diagonal boundary}
The two clauses quantify exactly over the image of g. In a virtual-class
theory allowing \(c=\{x\in D:x\notin g(x)\}\), with each
\(g(x)\) a class of members of D, c is not in \(g[D]\).
The diagonal conclusion needs this class-comprehension and membership
setup; it is not a theorem that every arbitrary abstract V exceeds D.
The source's extended numbers, projective ideal points and virtual classes
illustrate using ideal objects for regular notation, while any elimination
theorem depends on the particular theory chosen.
\end{hbtheorem}

\subsection{Truth concepts, partial terms and descriptions}

\begin{hbdefinition}{def:p08-sa-concepts}{Source concept types and local equality}
Write \(2=\{0,1\}\), with 1 true and 0 false. The source
\(\|\Phi\|\in2^I\) is a propositional concept, and
\(\|\Phi\|_i\) its value at i. Boolean connectives are evaluated
pointwise and remain two-valued. A term has an individual concept
\(\|\tau\|\in V^{(I)}\), where the parenthesized exponent denotes
\emph{partial} functions from I into V; \(V^I\) instead denotes
total functions. A virtual individual a and the rigid concept
\(\|\bar a\|\) are different typed objects.
For each fixed actual-domain set the source containment is
\(A_i^{(I)}\subseteq D^{(I)}\subseteq V^{(I)}\);
this does not make the varying \(A_j\) independent of j.
All displayed substitutions are capture-free, with bound variables
renamed when necessary.

For the source local equality, \(\|\tau=\sigma\|_i=1\) iff
both term values at i are defined and equal, or both are undefined.
If just one is undefined, equality is false. This is distinct from
the strict ordinary predicate convention below. The source intensional
equality is the formula
\(\tau\equiv\sigma\defeq\Box(\tau=\sigma)\), using its
\emph{universal} box. Here \(\equiv\) names equality of whole
individual concepts, not elementary equivalence of structures.

With the source description binder rendered \(\mathsf I\),
\(\|\mathsf I x\Phi(x)\|_i=a\) iff
\(\{b\in D:\|\Phi(\bar b)\|_i=1\}=\{a\}\).
Otherwise its value is undefined. Its dotted version uses \(A_i\)
instead of D. Consequently
\(\mathsf I_{\cdot}x\Phi(x)
\equiv\mathsf I x\exists_{\cdot}y(y=x\land\Phi(y))\).
Let \(\ast\in V^{(I)}\) be the totally undefined concept;
it is the value of \(\mathsf I x(x\ne x)\), not a new
arbitrarily selected individual of \(V\setminus D\)
\cite[\S III, pp.~152--155]{Scott1970}.
\end{hbdefinition}

\begin{hbtheorem}{thm:p08-sa-descriptions}{proved=annex:p8.sa.descriptions}{Every source description and equality case}
Local equality is reflexive, symmetric and transitive even with undefined
terms. Descriptions satisfy
\[
 \forall y\bigl(y=\mathsf I x\Phi(x)
       \leftrightarrow\forall x(x=y\leftrightarrow\Phi(x))\bigr)
\]
and
\[
 \neg\exists y(y=\mathsf I x\Phi(x))
       \to\mathsf I x\Phi(x)=\mathsf I x(x\ne x).
\]
These are pointwise statements, not an assertion that an improper
description is undefined at every other index. Equality of rigid constants
permits replacement in arbitrary capture-free formula contexts, and
\(\forall x\forall y((x=y\land\Phi(x))\to\Phi(y))\)
is valid. Local equality of arbitrary closed terms does not permit that
replacement; full concept equality \(\tau\equiv\sigma\) does.
\end{hbtheorem}

\subsection{Intensional predicates and propositional operations}

\begin{hbdefinition}{def:p08-sa-operators}{The source's operation types}
A general binary predicate has a total interpretation
\(\|R\|:V^{(I)}\times V^{(I)}\to2^I\), with
\(\|\tau R\sigma\|=\|R\|(\|\tau\|,\|\sigma\|)\).
An n-place predicate has the analogous n concept arguments.
It may inspect more than the two current extensions.
For an ordinary relation \(R_0\subseteq V\times V\), its barred
predicate is true at i exactly when both extensions are defined and
\((\|\tau\|_i,\|\sigma\|_i)\in R_0\).
A varying relation \(i\mapsto R_i\) gives the source tilded
predicate by the same strict clause at i. In particular
\(\|\widetilde A\tau\|_i=1\) iff the term is defined there
with value in \(A_i\).

The source general monadic propositional operation, rendered
\(\boxed{\mathscr P}\), has type \(2^I\to2^I\).
Its universal box sends an argument concept to the constant truth
concept precisely when that argument is 1 at every index. Its relational
operator \(\boxed\rho\), for \(\rho\subseteq I\times I\), is true
at i exactly when the argument is 1 at every j with \(i\rho j\).
If I has a composition \(\circ\), the earlier source example
tests the argument at all \(i\circ j\); this is the relational special case
\(i\rho k\) iff \(k=i\circ j\) for some j.
The source motivates this case with transformations satisfying
\(t_j(t_i(a))=t_{i\circ j}(a)\).

For \(I=\mathbb R\), the source progressive operator
\(\boxed{\leftrightarrow}\) is true at i when some open interval
around i is contained in the truth set of its argument. More generally
its arbitrary-family form is
\[
 \|\boxed{\mathscr P}\Phi\|_i=1
 \quad\Longleftrightarrow\quad
 \{j\in I:\|\Phi\|_j=1\}\in\mathscr P_i,
 \qquad \mathscr P_i\subseteq\mathcal P(I).
\]
This is Scott's published neighborhood clause, not an imported
axiomatization \cite[\S IV, pp.~155--161]{Scott1970}.
\end{hbdefinition}

\begin{hbexample}{def:p08-sa-pullback}{Separating varying predicates from index operations}
For a fixed relation \(R_0\subseteq V\times V\) and transformations
\(t_i:V\to V\), Scott proposes keeping individuals fixed while
setting \(aR_i b\) iff \((t_i(a),t_i(b))\in R_0\).
This preserves exactly the indicated atomic relation truth test;
it does not imply that the maps are bijective or preserve the whole
interpretation's equality and quantified-domain structure. The varying
predicate family and the chosen index operation are separately specified.
\end{hbexample}

\begin{hbtheorem}{thm:p08-sa-extensionality}{proved=annex:p8.sa.extensionality}{Pointwise extension versus composition of intensions}
Strict barred and tilded predicates preserve local equality of their
arguments; not every locally extensional predicate comes from an ordinary
relation, since the source equality makes two undefined arguments equal
where every strict binary relation is false. Full concept equality is
preserved by every general predicate and operation. For closed formulas,
\(\Box(\Phi\leftrightarrow\Psi)\) permits replacement at any chosen
occurrences, including under general intensional operations. Equality
only at the current index need not do so.
\end{hbtheorem}

\begin{hbtheorem}{thm:p08-sa-universal-quantifiers}{proved=annex:p8.sa.universal-quantifiers}{Universal necessity and the two quantified domains}
The source universal box validates K, T, 4, 5 and its duality, and preserves
logical validity under necessitation. Its pointwise truth does not mean
validity over every interpretation. Fixed possible-domain quantification
commutes with this box:
\(\Box\forall x\Phi(x)\leftrightarrow\forall x\Box\Phi(x)\).
Neither implication between \(\Box\forall_{\cdot}x\Phi(x)\)
and \(\forall_{\cdot}x\Box\Phi(x)\) is generally valid.
Actual quantification itself is exactly the restriction
\(\forall_{\cdot}x\Phi(x)\leftrightarrow
\forall x(\widetilde A x\to\Phi(x))\).
\end{hbtheorem}

\begin{hbtheorem}{thm:p08-sa-neighborhood-bridge}{proved=annex:p8.sa.neighborhood-bridge}{The exact modern neighborhood comparison}
Every map \(F:2^I\to2^I\) is uniquely represented by the source
families \(\mathscr P_i=\{X\subseteq I:F(1_X)(i)=1\}\).
Taking W=I and \(N(i)=\mathscr P_i\) preserves formula truth in
the common propositional language. Equal argument truth sets therefore
have equal modal outputs, providing theorem-level replacement of
equivalents, not replacement on an undischarged local premise.

As an editorial comparison, M, C and N correspond respectively to
upward closure, binary intersection closure and inclusion of I in each
family, with two atoms for arbitrary-set converse witnesses. Their
combinations are the later EM, EC, EN, EMC and EMCN/K presentations
of Sections~\ref{sec:p08-n-b-006-03} and \ref{sec:p08-n-b-006-04}.
EN need not distribute; regular EMC need not have a unit. The source
universal box is the special choice \(\mathscr P_i=\{I\}\),
not every member of that general family of modal systems.
\end{hbtheorem}

\begin{hbtheorem}{thm:p08-sa-progressive}{proved=annex:p8.sa.progressive}{The real-line example and the exact representation obstruction}
The progressive operator is topological interior and validates the S4
schemas. No relation on the same real-line domain represents this operation
for \emph{all} subsets of \(\mathbb R\): its neighborhoods at i
exclude their singleton core \(\{i\}\). This does not prevent a
different relational model from representing the same modal logic.
Scott attributes the stronger exact-S4 result to McKinsey and Tarski;
that completeness argument is not presented in this chapter
\cite[p.~160]{Scott1970}.
\end{hbtheorem}

\begin{hbwarning}{def:p08-sa-logical-operator-boundary}{Permutation invariance is a proposal, not the printed four-operation classification}
Scott suggests independence from the structure of I, mentioning permutation
invariance, as a possible way to delimit logical propositional operations.
He mentions negation, double negation, universal box and its dual.
The chapter gives no exact criterion or classification proof. Commuting
with all permutations alone does not single out only the operations he
mentions: the constant-true and constant-false maps also commute. Thus
this remark is not adopted as a proved classification. The source's
case for using several coexisting modal operators concerns the intended
coordinates and applications, not a mathematical demand that one
particular operator inventory must always be present.
\end{hbwarning}

\subsection{Binding operations and the source proof boundary}

\begin{hbdefinition}{def:p08-sa-binding}{The source general term/formula binder}
For the source placeholder \(\$x[\tau(x),\Phi(x)]\), all
bound x still range over D. Define whole-concept argument functions
\(f:D\to V^{(I)}\), \(f(a)=\|\tau(\bar a)\|\), and
\(F:D\to2^I\), \(F(a)=\|\Phi(\bar a)\|\).
For a term maker,
\[
 \|\$\|:(V^{(I)})^D\times(2^I)^D\to V^{(I)},\qquad
 \|\$x[\tau(x),\Phi(x)]\|=\|\$\|(f,F).
\]
A formula maker replaces the codomain by \(2^I\).
The source example \(\mathsf I y\forall x\Box(\Phi(x)\to y=\tau(x))\)
is a compound term of this format, but no theorem reducing all general
binders to descriptions and quantifiers is supplied. An interpretation
\(\mathfrak A\) comprises the domains and indices, actual-domain
family, rigid constants, logical clauses and type-correct nonlogical
operations. Logical validity quantifies over every such interpretation,
every free-variable assignment into D, and every index in its I
\cite[\S V, pp.~161--165]{Scott1970}.
\end{hbdefinition}

\begin{hbtheorem}{thm:p08-sa-binder-replacement}{proved=annex:p8.sa.binder-replacement}{Replacement through the source binders}
If \(\forall x(\tau(x)\equiv\sigma(x))\) and
\(\forall x\Box(\Phi(x)\leftrightarrow\Psi(x))\) hold,
then the corresponding binder outputs have the same full concept.
For term makers the conclusion is
\(\$x[\tau(x),\Phi(x)]\equiv\$x[\sigma(x),\Psi(x)]\);
for formula makers use the universal boxed biconditional.
Merely pointwise equal argument extensions do not suffice for a general
intensional binder.
\end{hbtheorem}

\begin{proofboundary}{def:p08-sa-completeness-boundary}{The actual 1970 axiomatization claim}
On p.~164 Scott reports that the valid first-order formulas are
axiomatizable by predicate logic, S5 for the universal box and replacement
of necessary equalities and biconditionals. He directs the Henkin argument
to planned joint work with Kaplan and notes Thomason's related fragments.
The chapter gives neither a complete formal rule inventory nor that
completeness proof. This is a \emph{reported source boundary}, not a
proved neighborhood completeness theorem. Higher-type extensions and
modal set theory are also indicated without a completed construction here.
The selected E/EM/EC/EN/EMC proofs are supplied under their actual later
sources, not retrospectively credited to Scott \cite[pp.~164--165]{Scott1970}.
\end{proofboundary}

\subsection{Incidence, equality and the published postscript}

\begin{hbdefinition}{def:p08-sa-incidence}{Individuals as states and the second function layer}
In the source's initial state representation, \(V\subseteq S^I\):
each individual a is a \emph{total} state history \(i\mapsto a_i\).
A term concept remains the different, possibly partial function
\(\|\tau\|:I\rightharpoonup V\). At i its denoted individual
\(\|\tau\|_i\) is itself a whole history; its current state is
\((\|\tau\|_i)_i\). Incidence \(\tau\approx\sigma\) holds
at i when these two states are defined and equal, or both term
values are undefined. Thus incidence, local equality of denoted
individuals and whole-concept equality are three separate relations.
With a restricted fixed property family, the state assignment may
instead fail to separate individuals over all indices; that is a
different interpretation hypothesis from identifying \(V\subseteq S^I\)
\cite[\S VI, pp.~165--169]{Scott1970}.
\end{hbdefinition}

\begin{hbtheorem}{thm:p08-sa-incidence}{proved=annex:p8.sa.incidence}{Conditional constant reconstruction and every surviving implication}
Under the injective history representation, rigid constants satisfy
\(\bar a=\bar b\leftrightarrow\Box(\bar a\approx\bar b)\).
For general closed terms the always valid arrows are
\[
 \tau\equiv\sigma\ \Longrightarrow\ \Box(\tau=\sigma)
 \ \Longrightarrow\ \begin{cases}\tau=\sigma,\\
                              \Box(\tau\approx\sigma),\end{cases}
 \qquad
 \tau=\sigma\Longrightarrow\tau\approx\sigma,
 \quad\Box(\tau\approx\sigma)\Longrightarrow\tau\approx\sigma.
\]
The first implication is an equivalence by definition.
Neither \(\tau=\sigma\) nor \(\Box(\tau\approx\sigma)\)
generally implies the other. Universal incidence does not imply whole
concept equality, even if all individuals have distinct complete histories.
For a nonseparating property family, even the constant reconstruction
needs an additional separation assumption. The annex gives explicit
finite histories rather than relying on an unresolved illustrative story.
\end{hbtheorem}

\begin{hbdefinition}{def:p08-sa-counterparts}{The source counterpart comparison}
For this separate construction let \(D=\bigcup_{i\in I}A_i\)
with pairwise disjoint actual domains. At i, write \(C_i(a,b)\)
for a being a current counterpart of b. Require
\(C_i(a,b)\Rightarrow a\in A_i\),
\(C_i(a,b)\land b\in A_i\Rightarrow a=b\), and
\(a\in A_i\Rightarrow C_i(a,a)\).
These are the truth conditions of the source tilded counterpart
predicate \(\widetilde C\). Its reading of counterpart necessity
for \(\Phi(x)\) is \(\forall y(y\widetilde C x\to\Phi(y))\).
The associated incidence is
\[
 \tau\approx\sigma\defeq\forall_{\cdot}x
       (x\widetilde C\tau\leftrightarrow x\widetilde C\sigma).
\]
This counterpart-profile incidence is an alternative construction,
not a claim that it always equals the state-history incidence above
\cite[pp.~169--171]{Scott1970}.
\end{hbdefinition}

\begin{hbtheorem}{thm:p08-sa-counterparts}{proved=annex:p8.sa.counterparts}{The source counterpart implication and its failed converse}
At each index, equality of counterpart profiles is an equivalence
relation on D. Moreover
\(\forall x\forall y((\widetilde A x\land x\approx y)
\to x\widetilde C y)\) is valid. Its converse is not:
one current counterpart can belong to two unequal profiles. Thus the
incidence relation does not recover the whole counterpart relation.
\end{hbtheorem}

\paragraph{The postscript belongs to the source's conclusions.}
Scott's December 1969 postscript treats the proposed framework as
unfinished philosophical analysis, not a final account of cross-world
identification. Kaplan distinguishes individuals from their concepts and
criticizes naming concept equality ``identity''. His highway examples
separate local merging from identical complete routes, and local equality
of descriptions from universal incidence. Montague's organism/molecules
comparison makes shared current state compatible with distinct individuals.
Their objections also challenge the choice of arbitrary intensional
predicate interpretations. The postscript distinguishes strict modal
logic, pragmatics with nonbinding propositional operators, extended
pragmatics with binders, and higher-order variants; it does not collapse
them into one theorem package \cite[pp.~171--173]{Scott1970}.
The mathematical distinctions and counterexamples above retain those
qualifications. They do not endorse the source's abandoned defense of
an unrestricted constant-incidence/equality biconditional.
